Engineers who depend on a slow simulation can afford only a limited number of runs, and must decide which fast stand-in model to train and whether more runs are worth the cost. We present SurroSelect, an open tool that chooses a stand-in from a pool of classical models by cross-validation, reports which models are tied, checks whether the POD basis limits accuracy, and forecasts the error for more runs. We tested it against criteria committed before the evaluation runs, first on \NTpProblems{} problems solved by the third-party finite-element library \texttt{scikit-fem} and then, with the tool frozen, on \NTpProblems{} further ones; we also compared it with a tuned neural network and ran it on two real measured datasets, and a \NAllProblems-problem study describes where each model family wins. On the further problems selection was within a median factor \SvSoneMedian{} of the oracle (the best pool member by test error), and always choosing one fixed pool member does almost as well. The tool was again worse than the benchmark's own selection among methods that include networks, in particular on wave problems, and the error forecast for four times the budget was within a factor two in only \SvpBoneShare{} of \SvpBoneN{} cases over all sixteen problems, usually too optimistic. Adding a tuned network to the pool helped on one of the two wave problems, at several times the cost. For a solver taking \SvETsim{} seconds per run, the surrogate was cheaper than the solver after about \SvEBreakTwofiftysix{} predictions. All problems are simulated and their solvers are fast; the evidence for much slower simulations is missing.
